\section{Integration with the BX3 Framework}
\label{sec:rs-integration}

Recursive Spawning does not operate as an independent mechanism. It is embedded
within and sustained by the BX3 three-layer architecture, which provides the
structural enforcement infrastructure that makes the spawn tree's accountability
guarantees unconditional.  The three layers — \purpose{}, \bounds{}, and
\fact{} — serve distinct but mutually reinforcing roles in the Recursive
Spawning pillar, each contributing an irreducible component to the overall
guarantee.  This section specifies how each layer participates in the Recursive
Spawning mechanism and how layer interactions enforce the invariants required
for drift prevention, chain integrity, and termination.

% ---------------------------------------------------------------
\subsection{Purpose Layer as Accountability Anchor}
\label{sec:purpose-layer-anchor}

The \purpose{} layer is the apex of the BX3 accountability architecture.  It
carries the human principal's intent and serves as the exclusive terminus for any
exception state that exceeds machine-actor provisioned authority.  In the
Recursive Spawning pillar, the \purpose{} layer plays two distinct roles that are
critical to the correctness of the spawn tree.

First, the \purpose{} layer defines the operating range $R(n)$ for every node
$n$ at provisioning.  This range is not merely a policy parameter — it is the
binding constraint that the \bounds{} layer enforces at every decision point.
The \purpose{} layer encodes the principal's intent as a structured scope that
constrains all downstream reasoning and action.  When a node spawns a child, the
\purpose{} layer's definition of $R(n)$ propagates downward as the upper bound
on $R(\text{child})$: the child cannot receive authority that its parent does not
possess.

Second, and more critically, the \purpose{} layer is the structural terminus for
the escalation chain.  Every node $n$ is assigned a human anchor $h(n)$ at
provisioning — a specific human principal who is the ultimate accountable party
for all actions taken under $R(n)$.  This assignment is a structural property of
the \purpose{} layer, not a configurable parameter.  The \purpose{} layer
carries the judgment required to resolve exceptions that exceed $R(n)$, and no
machine actor possesses the judgment capacity to substitute for it.  This is not
a policy choice — it is a structural consequence of what it means to carry
intent: an agent operating under delegated authority cannot possess the
principal's intent in sufficient completeness to serve as a substitute terminus
for unresolved exceptions.

The non-delegability of the \purpose{} layer's terminus function is the
architectural linchpin of the Recursive Spawning pillar.  It ensures that the
escalation chain from any node terminates at a human who is empowered to make
binding decisions, and that no machine actor can absorb or redirect this chain.
The \purpose{} layer is therefore the accountability anchor: the fixed point
to which all escalation paths converge, and the layer from which all legitimate
authority flows.

% ---------------------------------------------------------------
\subsection{Bounds Engine Depth Enforcement}
\label{sec:bounds-engine-depth}

The \bounds{} layer carries bounded reasoning and constrained execution within
$R(n)$.  In the Recursive Spawning pillar, the \bounds{} layer enforces the
depth constraints that prevent unbounded proliferation of the spawn tree.  This
enforcement is structural, not advisory.

At every spawn event, the \bounds{} layer evaluates whether the proposed child
node satisfies the hierarchy finiteness constraint: its depth $d(\text{child})$
must satisfy $d(\text{child}) \leq D_{\max}$, where $D_{\max}$ is the maximum
spawn depth established by the framework parameters.  This check is performed by
the \bounds{} layer's spawn state machine before the \texttt{child-provisioned}
event is emitted.  If the depth check fails, the state machine transitions to
\texttt{SPAWN\_LOCKED} and emits the $e_{\text{depth-exceeded}}$ event.  The
spawn event is rejected; no child node is created.

The \bounds{} layer additionally enforces the delegation scope inheritance
constraint.  Before any spawn event is committed, the \bounds{} layer evaluates
whether $R(\text{child}) \subset R(\text{parent})$ — whether the child's
operating range is a proper subset of the parent's.  This check prevents
authority expansion through recursive delegation: a parent cannot spawn a child
with greater authority than the parent itself possesses.  If the check fails, the
$e_{\text{scope-violation}}$ event is emitted, the spawn state machine
transitions to \texttt{SPAWN\_LOCKED}, and the Bailout Protocol is triggered.
The attempted expansion is both prevented and recorded.

These depth enforcement mechanisms are enforced by the \bounds{} layer's state
machine transition logic, which is itself subject to the \fact{} layer's
pre-commit hook.  The \bounds{} layer cannot bypass its own depth checks because
the checks are encoded in the state machine's transition guard conditions, not
in discretionary policy code.  The state machine will not transition to
\texttt{SPAWN\_ACTIVE} unless all guard conditions are satisfied.  This makes
depth enforcement a structural property of the spawn state machine, not an
operational one.

The $T_{\text{bounds}}$ timeout parameter is also enforced at the \bounds{}
layer, ensuring that bounded reasoning does not become unbounded deliberation.
If a node's reasoning process exceeds $T_{\text{bounds}}$, the \bounds{} layer
triggers the Bailout Protocol, preventing the node from continuing to act under
delegated authority without human review.  The timeout is a depth-equivalent
enforcement mechanism: it limits the operational duration of reasoning within
$R(n)$, complementing the depth limit on the spawn tree's structural nesting.

% ---------------------------------------------------------------
\subsection{Fact Layer Forensic Ledger}
\label{sec:fact-layer-fledger}

The \fact{} layer carries deterministic enforcement and forensic auditability.
It is the fail-safe anchor of the Recursive Spawning pillar: the only layer
that may issue a definitive halt to autonomous action, and the layer that
maintains the append-only Forensic Ledger that makes the spawn tree
reconstructable and tamper-evident.

The Forensic Ledger records every spawn event, every state transition of the
spawn state machine, every exception event, and every human directive issued
by any $h(n)$.  Each entry is a 6-tuple that includes the node identifier, the
logical timestamp, the event type, the full diagnostic context, and the SHA-256
hash of the preceding entry.  The cryptographic chaining ensures that no entry
can be modified, deleted, or inserted without breaking the chain — a condition
detectable by any auditor without access to the ledger's contents.

For Recursive Spawning specifically, the Forensic Ledger records the following
event types:

\begin{itemize}
\item \texttt{child-provisioned}$(p, c, R(c), h(c))$: records that parent $p$
  spawned child $c$ with operating range $R(c)$ and human anchor $h(c)$.
  This entry encodes the parent pointer $\alpha(c) = p$ and the human anchor
  $h(c)$ for the child node.

\item \texttt{spawn-locked}$(p, c, \text{reason})$: records that a spawn event
  was rejected, including the reason for rejection (depth exceeded, scope
  violation, or authorization failure).

\item \texttt{bailout-triggered}$(n, e)$: records that the Bailout Protocol was
  triggered at node $n$ due to exception $e$.

\item \texttt{parent-pointer-write-denied}$(v, \alpha_{\text{old}},
  \alpha_{\text{new}})$: records that a machine actor attempted to modify the
  parent pointer of node $v$ and was rejected by the \fact{} layer's pre-commit
  hook.
\end{itemize}

The Forensic Ledger is the sole source of truth for the runtime state of the
Recursive Spawning pillar.  The \fact{} layer's pre-commit hooks for the spawn
state machine query the Forensic Ledger to determine the current state —
specifically, the current depth $d(v)$ of any node $v$ and the current
$h(\cdot)$ assignments — before committing any state change.  This means that
the state machine's guard conditions are evaluated against the ledger's record,
not against a mutable in-memory state.  The ledger is the authoritative
checkpoint for all state machine decisions.

The Forensic Ledger's completeness invariant (Invariant
\ref{inv:forensic-completeness}) ensures that every spawn event is recorded
before subsequent actions are permitted.  This ordering constraint — record
before act — is enforced by the \fact{} layer's pre-commit hook.  No spawn
event can be committed without a corresponding \texttt{child-provisioned} entry
in the ledger.  This eliminates the possibility of a spawn event that is not
reflected in the audit record, ensuring that the chain from any node is fully
reconstructable from the ledger at any future point in time.

% ---------------------------------------------------------------
\subsection{Three-Layer Joint Enforcement}
\label{sec:three-layer-joint}

The correctness of Recursive Spawning depends on the conjunction of all three
layers acting in concert.  No single layer is sufficient, and no layer can
substitute for another.

The \purpose{} layer provides the accountability anchor: the human principal who
is the ultimate terminus for all escalation paths and the source of all
legitimate authority.  Without the \purpose{} layer, there would be no fixed
point to terminate the chain, and no source of intent to propagate downward
through the spawn tree.

The \bounds{} layer provides the depth enforcement: the state machine guard
conditions that prevent unbounded spawning, authority expansion, and unbounded
reasoning.  Without the \bounds{} layer, the spawn tree could grow without
structural limit and the escalation chain could be redirected by a machine actor
with sufficient procedural sophistication.

The \fact{} layer provides the forensic ledger: the append-only record that
makes the spawn tree tamper-evident and reconstructable, and the pre-commit
hooks that enforce the state machine's guard conditions at runtime.  Without
the \fact{} layer, the spawn tree's state would be mutable, the audit record
would be suppressible, and the state machine's guard conditions could be
bypassed.

The joint enforcement is formalized in the three-layer interaction constraints
(IC-1 through IC-4) specified in Section~\ref{sec:layer-constraints}.  These
constraints are structural — enforced by the \fact{} layer's pre-commit hook,
not by convention or policy.  They ensure that each layer performs its
designated function, that no layer can substitute for another, and that the
accountability chain from any node to its human anchor is preserved regardless
of what any machine actor attempts to do at runtime.
